\section*{第 \ref{ch:b2:charges-currents-conduction} 章——电荷、电流与导电}

\begin{solution}{exo:b2:charges-currents-conduction:1}
$j = 16/2.5 \times 10^{-6} = \qty{6.4e6}{A/m^2}$; $v = j/ne = \qty{0.47}{mm/s}$; \qty{20}{m}
所需时间 $4.3 \times 10^4$ s，即十二小时。$E = j/\gamma = \qty{0.11}{V/m}$；$U = \qty{2.1}{V}$；
$P = UI = \qty{34}{W}$。
\end{solution}

\begin{solution}{exo:b2:charges-currents-conduction:2}
闪电 $3 \times 10^4/\pi10^{-4} = \qty{1e8}{A/m^2}$；光束 $\qty{4e3}{A/m^2}$；轴突
$10^{-9}/1.3 \times 10^{-11} = \qty{80}{A/m^2}$；环电流 \qty{1e-8}{A/m^2}。只有
闪电超过导线的 $6 \times 10^6$。
\end{solution}

\begin{solution}{exo:b2:charges-currents-conduction:3}
$R = \rho_eL/S$：$0.68$、$1.1$、$4.0$、\qty{44}{\ohm}。铝：$2.5 \times 2.8/1.7 =
\qty{4.1}{mm^2}$，质量之比 $(4.1 \times 2700)/(2.5 \times 8900) = 0.50$——质量减半。
加热器：$R = 230^2/1000 = \qty{53}{\ohm}$，$S = \qty{0.196}{mm^2}$，$L = RS/\rho_e = \qty{9.4}{m}$。
\end{solution}

\begin{solution}{exo:b2:charges-currents-conduction:4}
$\tau = m\gamma/ne^2$：铜 \qty{2.5e-14}{s}，$\mu = e\tau/m = \qty{4.4e-3}{m^2/(V.s)}$，
$\ell = \qty{40}{nm}$；银 \qty{3.8e-14}{s}，\qty{6.7e-3}{m^2/(V.s)}，\qty{61}{nm}。
硅：$\gamma = ne\mu = 10^{22} \times 1.6 \times 10^{-19} \times 0.14 = \qty{220}{S/m}$；$v =
\mu E = \qty{140}{m/s}$。
\end{solution}

\begin{solution}{exo:b2:charges-currents-conduction:5}
(a) 电流流入而无流出：$\dd Q/\dd t = I$。(b) 电荷累积：$\partial_t
\rho \ne 0$。(c) $\vect j = (I/4\pi r^2)\vect e_r$：$\operatorname{div}\vect j = \frac1{r^2}
\partial_r(I/4\pi) = 0$。(d) 仅在电极处成立，此处电流
由导线注入——是通量的一个源。
\end{solution}

\begin{solution}{exo:b2:charges-currents-conduction:6}
(a) $j = I/2\pi rL$，$E = j/\gamma$，沿径向。(b) $U = \int_a^bE\,\dd r = I\ln(b/a)/
2\pi\gamma L$：$R = \ln3.5/(2\pi \times 10^{-13} \times 10^3) = \qty{2e9}{\ohm}$ 每千米。
(c) \qty{0.5}{\micro A}，\qty{0.5}{mW}。(d) $RC = \varepsilon_0\varepsilon_r/\gamma = \qty{200}{s}$（取
$\varepsilon_r = 2.3$）：带电并与外界隔离的电缆，通过其自身的绝缘体放电，
放电时间尺度即该绝缘体的弛豫时间。
\end{solution}

\begin{solution}{exo:b2:charges-currents-conduction:7}
(a) $\omega\tau$：$8 \times 10^{-12}$、$1.6 \times 10^{-3}$、$4.7$、$79$。(b) $|\underline\gamma|/\gamma_0 =
1/\sqrt{1 + \omega^2\tau^2}$：$1$、$1$、$0.21$、$0.013$；相位 $-\arctan\omega\tau$：$0$、
$-\qty{0.09}{\degree}$、$-\qty{78}{\degree}$、$-\qty{89}{\degree}$。(c) $\omega\tau \le 0.14$：$f \le
\qty{0.9}{THz}$。(d) $\underline\gamma \approx ne^2/\iu m\omega$：$\underline{\vect j}$ 相对
$\underline{\vect E}$ 落后 \qty{90}{\degree}，$\langle\vect j\cdot\vect E\rangle = 0$：没有焦耳
热——电子像在等离子体中那样自由振荡。
\end{solution}

\begin{solution}{exo:b2:charges-currents-conduction:8}
(a) $S = \qty{3.1e-8}{m^2}$，$j = \qty{3.2e8}{A/m^2}$，$p = j^2/\gamma = \qty{1.7e9}{W/m^3}$。
(b) 单位体积熔化所需能量 $\rho(c\Delta T + L_f) = 8900(4.1 \times 10^5 + 2.05 \times
10^5) = \qty{5.5e9}{J/m^3}$：\qty{3.2}{s}。(c) 快九倍：\qty{0.36}{s}。(d) 额定电流
是这样一个电流：其焦耳热恰好被空气带走，
使温度保持在熔点以下；沙子用来熄灭电弧，
否则导线熔断后，电弧会跨越间隙继续导电。
\end{solution}

\begin{solution}{exo:b2:charges-currents-conduction:9}
(a) $\partial_t\rho = -\operatorname{div}\vect j = -\gamma\operatorname{div}\vect E = -\gamma\rho/
\varepsilon_0$。(b) $\varepsilon_0/\gamma$：\qty{1.5e-19}{s}、\qty{1.8e-12}{s}、\qty{1.8e-10}{s}、
\qty{9}{s}、$9 \times 10^4$ s。(c) $f \lesssim 1/2\pi\tau_r$：$10^{18}$ Hz（\omterm{prop:b2:charges-currents-conduction:drude}{德鲁德模型}
早在此之前就已失效）、\qty{90}{GHz}、\qty{0.9}{GHz}、\qty{0.02}{Hz}、\qty{2e-6}{Hz}。(d)
玻璃棒若通过其体积导电，一分钟内就会失去电荷；
真实干燥玻璃的电导率更接近 $10^{-14}$ S/m（需数小时），通常是表面
湿度使其放电——在数量
级上相符。
\end{solution}

\begin{solution}{exo:b2:charges-currents-conduction:10}
(a) $\vect j = (I/2\pi r^2)\vect e_r$，$E = I/2\pi\gamma r^2$，$V = I/2\pi\gamma r$。(b) $R = V(a)
/I = 1/2\pi\gamma a = \qty{32}{\ohm}$。(c) $V_{\text{棒}} = \qty{640}{kV}$；$V(10) = \qty{32}{kV}$，
$V(10.7) = \qty{30}{kV}$：两脚之间为 \qty{2.1}{kV}。(d) 前
腿与后腿沿半径方向相距 \qty{1.5}{m}：$4$--\qty{5}{kV}，且
电流路径穿过心脏。
\end{solution}

\begin{solution}{exo:b2:charges-currents-conduction:11}
(a) $E_H = jB/nq$，其中 $j = I/wt$：$V_H = E_Hw = IB/nqt$。(b) $1/(8.5 \times 10^{28}
\times 1.6 \times 10^{-19} \times 10^{-4}) = \qty{0.7}{\micro V}$。(c) \qty{6.3}{mV}：$V_H \propto 1/n$，
而\omterm{rem:b2:charges-currents-conduction:materials}{半导体}的载流子数少 $10^6$ 倍。(d) $B = V_Hnqt/I =
10^{-5} \times 10^{22} \times 1.6 \times 10^{-19} \times 2 \times 10^{-5}/10^{-3} = \qty{0.32}{T}$；由
$V_H$ 的符号可得载流子的符号（对已知的 $\vect B$），或由
$\vect B$ 的方向（对已知的传感器）。
\end{solution}

\begin{solution}{exo:b2:charges-currents-conduction:12}
(a) $\vect j = ne\mu_+\vect E + (-ne)(-\mu_-\vect E)$：两种粒子
沿相反方向漂移，却沿同一方向输运电流。(b) $n =
0.5 \times 10^3 \times 6.02 \times 10^{23} = \qty{3e26}{m^{-3}}$：$\gamma = 3 \times 10^{26} \times 1.6 \times 10^{-19}
\times 1.31 \times 10^{-7} = \qty{6.3}{S/m}$（离子在此浓度下相互
阻碍，实测为 \qty{5}{S/m}）。(c) \qty{0.52}{\micro m/s} 与 \qty{0.79}{\micro m/s}：
Cl$^-$ 携带 $60\%$ 的电流。(d) $R = L/\gamma S = \qty{200}{\ohm}$，$U = \qty{200}{V}$，$P =
\qty{200}{W}$ 注入 \qty{10}{g} 水中：\qty{2}{s} 内升温 \qty{10}{K}——因此电泳槽
需配冷却极板。
\end{solution}

\begin{solution}{pb:b2:charges-currents-conduction:1}
\textbf{1.} $n = \rho N_A/M = \qty{8.4e28}{m^{-3}}$。

\textbf{2.} $j = \qty{6.4e6}{A/m^2}$，$v = j/ne = \qty{0.47}{mm/s}$；每米需 \qty{35}{min}。

\textbf{3.} $\tau = m\gamma/ne^2 = \qty{2.5e-14}{s}$，$\mu = \qty{4.4e-3}{m^2/(V.s)}$，$\ell =
\qty{40}{nm}$，$150$ 个间距：电子与缺陷及
热振动发生碰撞，而非与原子碰撞。

\textbf{4.} $\rho_e \propto T$ 给出 $\alpha = 1/293 = \qty{3.4e-3}{K^{-1}}$，与
实测值 $4 \times 10^{-3}$ 接近。

\textbf{5.} $E = \qty{0.11}{V/m}$，$U = \qty{0.11}{V}$；$\tfrac12mv^2 = \qty{1e-37}{J}$，与
$k_BT = \qty{4e-21}{J}$ 相比，电子气体并未察觉到电流。

\textbf{6.} $I/e = \qty{1e20}{}$，即每秒电子数。

\textbf{7.} $Q = neSL = \qty{3.4e4}{C}$；$QE = \qty{3.6}{kN}$，与
\qty{0.2}{N} 的重力相比。电场以
相反的力拉动带正电的晶格；电子通过碰撞把动量交给晶格：
整根导线整体上不受任何作用。

\textbf{8.} $p = j^2/\gamma = \qty{6.8e5}{W/m^3}$，\qty{1.7}{W/m}。

\textbf{9.} 每米表面积 $\pi d = \qty{5.6e-3}{m^2}$：$\Delta T = 1.7/(10 \times 5.6 \times
10^{-3}) = \qty{30}{K}$；电阻率升高 $12\%$，功率升至 \qty{1.9}{W}：
$\Delta T \approx \qty{35}{K}$。

\textbf{10.} $\Delta T = 1.7(1 + 0.004\Delta T)/(2 \times 5.6 \times 10^{-3})$：$\Delta T(1 - 0.6) = 150$，
$\Delta T \approx \qty{400}{K}$——接近热失控：埋地电缆必须
降额使用。

\textbf{11.} $p/\rho c = \qty{0.2}{K/s}$；在 \qty{160}{A} 下，升温速率为 \qty{20}{K/s}，
约一分钟内熔化（$\rho_e$ 升高时所需时间更短）；断路器在毫秒内断开。

\textbf{12.} $\omega\tau = 8 \times 10^{-12}$：欧姆定律不受影响。

\textbf{13.} $\varepsilon_0/\gamma = \qty{1.5e-19}{s}$：导线内部无电荷；导线中的电场
由其\emph{表面}上的电荷产生，该电荷面密度
沿导线变化。

\textbf{14.} $Q = I_0\tau_c = \qty{16}{mC}$，$V = Q/C = \qty{160}{V}$。

\textbf{15.} 导线一侧通有 $I$，间隙一侧无
传导电流：电荷不断增加，间隙中的电场也随之增强。

\textbf{16.} $I = 16\eu^{-0.5} = \qty{9.7}{A}$; $E = Q/\varepsilon_0A$, $\dd E/\dd t = I/\varepsilon_0A
= 9.7/(8.85 \times 10^{-12} \times 10^{-4}) = \qty{1.1e16}{V.m^{-1}.s^{-1}}$.

\textbf{17.} $\varepsilon_0A\,\dd E/\dd t = I$：位移电流。

\textbf{18.} $j = \gamma E = 10^{-14} \times 2.3 \times 10^5 = \qty{2e-9}{A/m^2}$, $10^{-14}$ A
通过截面的电流可忽略；在一毫米距离上击穿所需电压为 \qty{3}{kV}。

\textbf{19.} $R = 1000/(6 \times 10^7 \times 10^{-5}) = \qty{1.7}{\ohm}$; $\Delta U = \qty{67}{V}$; $P =
\qty{2.7}{kW}$，占所输送功率 \qty{9.2}{kW} 的 $29\%$。

\textbf{20.} $I = \qty{0.46}{A}$，$\Delta U = \qty{0.8}{V}$，\qty{0.35}{W}：采用高压输电时
损耗除以 $(U_2/U_1)^2$。

\textbf{21.} $10 \times 6/3.6 = \qty{17}{mm^2}$；每千米铜的质量为 \qty{89}{kg/km}，
对比 \qty{45}{kg/km} 的铝：架设在杆塔上更轻、更便宜；墙内
铜的截面较小，适用于线管和接线柱。

\textbf{22.} $V_H = IB/nqt = 5 \times 10^{-3} \times 0.1/(10^{22} \times 1.6 \times 10^{-19} \times 10^{-4}) =
\qty{3.1}{mV}$.

\textbf{23.} $+24\%$：\qty{2.1}{\ohm}，损失 \qty{3.3}{kW}。

\textbf{24.} $RI^2t = 1.7 \times 4 \times 10^8 \times 10^{-4} = \qty{67}{kJ}$ 注入 \qty{89}{kg} 的材料中，温升 \qty{2}{K}。
$j = \qty{2e9}{A/m^2}$，是家用数值的三百倍。

\textbf{25.} $n$（物质）与 $\tau$（碰撞）给出 $\gamma = ne^2\tau/m$；
$\gamma$ 联系 $j$ 与 $E$（欧姆定律）；$j$ 积分得 $I$，$E$ 积分得 $U$（
电路）；$j^2/\gamma$ 即为焦耳热。
\end{solution}
